\section{Метод решения}

Работа выполнена на C++17. Стандартные средства языка используются для строк,
разбора данных и управления временем жизни ресурсов, но процессы и каналы
создаются прямыми системными вызовами. \texttt{std::thread},
\texttt{std::filesystem} и другие высокоуровневые обёртки над ОС не
используются.

Родитель создаёт два однонаправленных канала. Первый передаёт строки с числами
дочерней программе, второй возвращает односимвольный статус. После каждой
команды родитель ожидает ответ, поэтому команды и статусы однозначно
соответствуют друг другу.

В Linux вызов \texttt{fork()} создаёт новый процесс. В дочерней ветви
\texttt{dup2()} переназначает читающий конец первого канала на стандартный
ввод, а пишущий конец второго канала --- на стандартный вывод. После этого
\texttt{execl()} заменяет образ процесса программой \texttt{child}. Родитель
использует \texttt{write()}, \texttt{read()} и \texttt{waitpid()}.

В Windows каналы создаются функцией \texttt{CreatePipe()}. Дескрипторы
передаются дочернему процессу как \texttt{STDIN} и \texttt{STDOUT} через
\texttt{STARTUPINFOA}; процесс создаётся функцией \texttt{CreateProcessA()}.
Передача данных выполняется через \texttt{WriteFile()} и \texttt{ReadFile()}.

Проверка простоты перебирает нечётные делители до квадратного корня. Условие
\texttt{divisor <= value / divisor} исключает переполнение. Значения 0 и 1
игнорируются, так как не являются ни простыми, ни составными.

\section{Описание программы}

\begin{itemize}
    \item \texttt{inc/parent.cpp} --- создание каналов и процесса,
    пользовательский ввод, отправка команд, чтение статусов;
    \item \texttt{inc/child.cpp} --- разбор чисел, проверка простоты,
    запись составных значений;
    \item \texttt{inc/common.cpp} --- общие функции разбора и проверки
    простоты;
    \item \texttt{inc/include/protocol.hpp} --- коды сообщений протокола;
    \item \texttt{CMakeLists.txt} --- кроссплатформенная сборка двух
    программ.
\end{itemize}

Дескрипторы Linux и Windows помещаются в собственные RAII-классы
\texttt{unique\_fd} и \texttt{unique\_handle}. Их деструкторы вызывают
\texttt{close()} или \texttt{CloseHandle()}, поэтому ресурс освобождается
при любом выходе из области видимости.